\section{Background and Motivation}
\label{sec:background}

\subsection{MoE Serving}
\MoE models replace dense feed-forward layers with a set of experts and a router.
For each token and layer, the router selects a small number of experts, typically top-1 or top-2.
The benefit is high parameter count at lower per-token compute; the cost is serving complexity.
Expert weights may be sharded across devices, moved between CPU and GPU, or dynamically rebalanced.
Late knowledge of expert choices creates stalls when the needed expert is not resident or when many tokens concentrate on the same expert.

Serving systems therefore rely on a mixture of expert parallelism, load balancing, prefetching, and placement.
Reactive policies observe recent routing and move experts accordingly.
These policies are model-agnostic but have limited visibility into upcoming requests.
\project asks whether agent workloads provide earlier request-level signals that complement router-time observation.

\subsection{Agent Workloads}
Agent runtimes build prompts from structured nodes.
A coding workflow may contain a planner, coder, tester, and critic.
A search workflow may contain a planner, searcher, and summarizer.
Each node has metadata before inference: role, phase, tool type, graph-node type, prompt-block types, ready time, dependency list, and criticality.
These fields are stable enough for scheduling and audit logs; \project reuses them as predictor keys.

The key distinction from semantic prediction is timing.
\project does not inspect generated text to predict future routing.
It uses pre-decode metadata that the scheduler already knows.
This makes the signal available early enough for expert prefetch and placement.

\subsection{Dense-Model Boundary}
Not every model has routed experts.
The requested end-to-end target, \Qwen at \texttt{/home/youwei/bzh/model/Qwen/Qwen2.5-7B-Instruct}, is a dense Qwen2 causal language model.
Its Hugging Face config contains \texttt{model\_type=qwen2}, architecture \texttt{Qwen2ForCausalLM}, and no expert-count fields.
For such models, \project's correct behavior is not to emulate a router.
Instead, the vLLM sidecar records metadata and explicitly returns \texttt{dense\_model\_no\_moe\_router}.
This negative case is part of the system contract and prevents misleading MoE claims.
